\documentclass[11pt, aspectratio=169]{beamer}
\usepackage[utf8]{inputenc}
\usepackage[T1]{fontenc}
\usepackage{amsmath, amssymb, amsfonts}
\usepackage{graphicx}
\usepackage{hyperref}
\usepackage{xcolor}

% Beamer theme settings
\usetheme{Madrid}
\usecolortheme{default}
\setbeamertemplate{navigation symbols}{}

\title{Action Chunking with Transformers (ACT)}
\subtitle{Data Pipeline, Training, and Inference}
\author{High-Level Overview}
\date{\today}

\begin{document}

\begin{frame}
    \titlepage
\end{frame}

\begin{frame}{Outline}
    \item Data preprocessing
    \item Training forward pass
    \item Loss and optimization
    \item Inference
\end{frame}

\section{Data Pipeline}

\begin{frame}{Introduction to ACT}
    \begin{itemize}
        \item \textbf{Action Chunking:} Predicts sequences of actions (chunks) rather than a single step, reducing compounding errors in fine manipulation.
        \item \textbf{Modular Design:} Easily tailored for symmetric or asymmetric setups (e.g., single robotic arm with 2-3 cameras).
        \item \textbf{CVAE Architecture:} Uses a Conditional Variational Autoencoder approach to capture the "style" of human demonstrations.
        \item \textbf{High-Level Flow:} Real-time vision + Proprioception (joint state) + Style variable $\rightarrow$ Action sequence.
    \end{itemize}
\end{frame}

\begin{frame}{Data Normalization}
    \begin{itemize}
        \item \textbf{Why Normalization?} Ensures numerical stability. Prevents motors with large ranges from overpowering fine gripper movements.
        \item \textbf{Vision Data:} 
        \begin{itemize}
            \item Scales pixel values to $[0, 1]$.
            \item Normalizes using distribution mean and standard deviation to center color channels and stabilize ResNet layers.
        \end{itemize}
        \item \textbf{Proprioception \& Actions:}
        \begin{itemize}
            \item Standardized using dataset statistics: $x_{norm} = \frac{x_{raw} - \mu}{\sigma}$.
            \item Keeps all joint dimensions on a similar scale. Process is reversed before sending commands to hardware.
        \end{itemize}
    \end{itemize}
\end{frame}

\begin{frame}{Processing Vision and State}
    \begin{itemize}
        \item \textbf{Vision Pipeline (Cameras):}
        \begin{itemize}
            \item Images pass through a \textbf{ResNet-18} to extract feature grids.
            \item Grids are flattened into tokens (e.g., 300 per camera).
            \item \textbf{2D Positional Embeddings} are added so the network remembers \textit{where} features appeared spatially.
            \item Tokens from all cameras are concatenated into a single visual sequence.
        \end{itemize}
        \item \textbf{Robot State (Proprioception):}
        \begin{itemize}
            \item The lower-dimensional joint state vector maps through a linear layer.
            \item Produces a \textbf{single state token} matching the Transformer's hidden dimension.
        \end{itemize}
    \end{itemize}
\end{frame}

\section{Training Forward Pass}

\begin{frame}{The CVAE Encoder (Capturing Style)}
    \begin{itemize}
        \item \textbf{Goal during training:} Learn the nuance of a human demonstration before the main Transformer handles the scene.
        \item \textbf{Inputs to the Style Encoder:} 
        \begin{itemize}
            \item Current physical state (no images, to save computing power).
            \item The exact \textbf{future actions} taken by the human.
        \end{itemize}
        \item \textbf{The [CLS] Token:} A designated "dummy" token that interacts bidirectionally with the whole sequence. It acts as a final summary bucket, discarding individual sequences to extract a single \textbf{Style Vector}.
    \end{itemize}
\end{frame}

\begin{frame}{Distributions and Sampling}
    \begin{itemize}
        \item Instead of a static tensor, the style vector represents a \textbf{Gaussian distribution}.
        \item The network predicts the statistical center (Mean $\mu$) and uncertainty (Variance $\sigma^2$).
        \item \textbf{The Reparameterization Trick:} To allow training (backpropagation), we sample the style $z$ using random noise: $z = \mu + \sigma \cdot \epsilon$
        \item This $z$ captures mathematical stochasticity and specific multi-modal aspects of the demonstration block.
    \end{itemize}
\end{frame}

\begin{frame}{The Policy Transformer (Action Generation)}
    \begin{itemize}
        \item \textbf{Policy Encoder:} Combines vision tokens, state token, and the sampled style token ($z$).
        \begin{itemize}
            \item Employs self-attention so these modalities can "talk" and produce a unified memory of the environment and intended style.
        \end{itemize}
        \item \textbf{Policy Decoder:} Queries the Encoder's memory using fixed position blocks.
        \begin{itemize}
            \item Generates a purely mathematical "sentence" representing the future path.
            \item A final linear layer crushes these high-dimensional tokens back directly into physical motor instructions (the Action Chunk).
        \end{itemize}
    \end{itemize}
\end{frame}

\section{Loss and Optimization}

\begin{frame}{The Dual Loss Objective}
    The network improves by minimizing two mathematical penalties simultaneously:
    \vspace{1em}
    \begin{enumerate}
        \item \textbf{L1 Reconstruction Loss:} Measures how closely the predicted chunk matches the human expert's ground-truth chunk.
        \begin{itemize}
            \item Uses L1 over Mean Squared Error (MSE) because physical robotics requires tight adherence to precise, fine-grained details without blurring.
        \end{itemize}
        \item \textbf{KL Divergence:} A penalty restricting how the "style" variable ($z$) behaves in multi-modal situations.
    \end{enumerate}
\end{frame}

\begin{frame}{The Multi-Modal Averaging Problem}
    \begin{itemize}
        \item \textbf{The Problem:} Humans demonstrate tasks with natural variation (e.g., grabbing a cup by the handle or by the top).
        \item Without a style variable, a basic network tries to mathematically average these conflicting demonstrations, resulting in poor execution (crashing into the cup).
        \item \textbf{The Solution ($z$):} The style variable $z$ gives the model an excuse to safely separate these approaches. Positive $z$ = Grab the handle; Negative $z$ = Grab the top.
    \end{itemize}
\end{frame}

\begin{frame}{Why the KL Divergence Prior?}
    \begin{itemize}
        \item We apply a massive penalty to force all natural style variations to cluster tightly around the mathematical origin $\mathcal{N}(0,I)$.
        \item By penalizing large values, the origin ($z=0$) becomes the \textbf{Canonical Baseline}---the absolute safest standard way to complete the task.
        \item It solves the issue for inference: when the human expert isn't there to provide a demonstration, setting $z=0$ cleanly breaks ties and prevents erratic behavior.
    \end{itemize}
\end{frame}

\begin{frame}{Update of Weights (Optimization Architecture)}
    \begin{itemize}
        \item The entire CVAE and Policy networks are optimized simultaneously. Gradients flow backwards directly from the Decoder outputs to the Style Encoder mapping.
        \item \textbf{AdamW Optimizer:} Preferred because it smoothly handles Transformer topologies.
        \begin{itemize}
            \item \textbf{Momentum:} Accelerates training in consistent downward slopes.
            \item \textbf{Adaptive Learning:} Individually scales learning rates up or down per weight, stabilizing deep internal layers.
            \item \textbf{Weight Decay:} A regularization technique that stops any single weight from ballooning, combatting overfitting to noise.
        \end{itemize}
    \end{itemize}
\end{frame}

\section{Inference}

\begin{frame}{The Simplified Inference Pass}
    \begin{itemize}
        \item In deployment, there is no human future data to interpret. The \textbf{Style Encoder is disabled and discarded}.
        \item \textbf{Hardcoding $z$:} We forcefully set $z = \vec{0}$, dropping the robot directly onto its safest, most robust core strategy.
        \item \textbf{Fast Loop:} The active pipeline distills down to merely reading cameras and proprioception, passing them to the Policy Encoder with $z=0$, and executing the chunk.
    \end{itemize}
\end{frame}

\begin{frame}{The Overlapping Chunk Problem}
    \begin{itemize}
        \item A single prediction provides a 100-step future action sequence.
        \item Playing 100 steps blindly to the end causes highly unreactive robotic motion---the robot wouldn't notice if something shifted.
        \item \textbf{High Frequency Queries:} ACT re-runs inference every few milliseconds.
        \item \textbf{The Conflict:} Because of constant replanning, the framework ends up with a queue of \textit{overlapping predictions competing for the exact same physical motor execution} at any given millisecond.
    \end{itemize}
\end{frame}

\begin{frame}{Temporal Ensembling}
    \begin{itemize}
        \item Overlapping chunk conflicts are solved using \textbf{Temporal Ensembling}.
        \item Instead of letting the newest chunk override older ones, they are mathematically blended together via \textbf{Exponential Weighting}.
        \item \textbf{Oldest Predictions (High Weights):} Provide a stable, committed trajectory core over time.
        \item \textbf{Newest Predictions (Low Weights):} Act as subtle "corrective nudges" accommodating immediate visual changes.
        \item \textbf{Result:} Smooth, continuous, remarkably reactive long-horizon policies.
    \end{itemize}
\end{frame}

\end{document}
